\chapter[Układy hiperboliczne i propagacja]{Równania hiperboliczne: układy symetryczno-hiperboliczne, energia i propagacja informacji}
\label{chap:pde21}

\section{Dlaczego potrzebujemy układów pierwszego rzędu?}
Rozdział~\ref{chap:pde-wave} dotyczył klasycznego równania falowego, przede wszystkim
jego jawnych formuł. Obecnie rozważamy układy o zmiennych współczynnikach,
bez dostępnej formuły rozwiązania. Ich centralną własnością będzie możliwość
kontrolowania energii bez tracenia pochodnych. Pozwoli to oddzielić
\emph{hiperboliczność symbolu} od \emph{poprawności zagadnienia Cauchy'ego}.

\begin{uwaga}[Dwie zasady energetyczne: dyfuzja i propagacja]
\label{rem:pde-v43-energy-types}
Oszacowanie paraboliczne z
twierdzenia~\ref{thm:pde20-existence} kontroluje także
całkę w czasie z normy przestrzennej $H_0^1$ dzięki
koercywności operatora dyfuzji. W układzie
symetryczno-hiperbolicznym energia kontroluje przede
wszystkim normę $L^2$ w chwili $t$, natomiast jej
lokalny bilans daje stożki zależności. Nie wolno
wywnioskować wygładzania przestrzennego z samej
nierówności energetycznej dla układu hiperbolicznego.
\end{uwaga}

Niech $d,m$ będą dodatnimi liczbami całkowitymi i $0<T<\infty$.
Dla kolumny $u=(u_1,\ldots,u_m)^\mathsf T$
rozważmy na $[0,T]\times\R^d$ równanie
\begin{equation}\label{eq:pde21-main}
 \partial_tu+\sum_{j=1}^d A_j(t,x)\partial_{x_j}u+B(t,x)u=f(t,x),
 \qquad u(0,x)=u_0(x).
\end{equation}
Macierze $A_j,B$ mają rozmiar $m\times m$ i rzeczywiste współczynniki.
Pracujemy nad ciałem liczb rzeczywistych. Przez $H$
oznaczamy przestrzeń Hilberta $L^2(\R^d;\R^m)$
z iloczynem skalarnym
$(v,w)_H=\int_{\R^d}v(x)^\mathsf Tw(x)\,\mathrm d\lambda_d(x)$
i normą $\|v\|_2^2=\int_{\R^d}|v(x)|^2\,\mathrm d\lambda_d(x)$.
Norma $|\cdot|$ w $\R^m$ jest normą euklidesową,
a $\|\cdot\|_{\mathrm{op}}$ jest indukowaną przez nią
normą operatorową macierzy. Przestrzenie funkcji czasu
o wartościach w przestrzeniach Sobolewa i w $H$
rozumiemy w sensie Bochnera.

\begin{definicja}[Układ symetryczno-hiperboliczny]\label{def:pde21-sym}
Układ~\eqref{eq:pde21-main} nazywamy \emph{symetryczno-hiperbolicznym},
jeżeli $A_j(t,x)=A_j(t,x)^\mathsf T$ dla wszystkich $j$ i wszystkich
$(t,x)$ (lub prawie wszędzie, gdy pracujemy ze współczynnikami mierzalnymi).
Jego symbolem głównym (z pominięciem wspólnego
czynnika $\mathrm i$ w konwencji Fouriera) jest
$P(t,x;\tau,\xi)=\tau I_m+\sum_{j=1}^d\xi_jA_j(t,x)$.
Dla współczynników mierzalnych rozpatrujemy symbol
dla prawie każdego $(t,x)$.
Kowektor $(\tau,\xi)\ne0$ jest charakterystyczny, jeżeli
$\det P(t,x;\tau,\xi)=0$.
\end{definicja}

\begin{stwierdzenie}[Rzeczywiste prędkości charakterystyczne]
\label{prop:pde21-real}
W każdym punkcie $(t,x)$, w którym wszystkie macierze
$A_j$ są symetryczne, dla każdego
$\xi\in\R^d\setminus\{0\}$ wszystkie pierwiastki wielomianu
$\tau\mapsto\det(\tau I_m+\sum_j\xi_jA_j)$ są rzeczywiste.
Płaszczyzna $t=\mathrm{const}$ jest niecharakterystyczna
niezależnie od symetrii $A_j$.
\end{stwierdzenie}
\begin{proof}
Macierz $A(\xi)=\sum_j\xi_jA_j$ jest rzeczywista i symetryczna.
Z twierdzenia spektralnego istnieje baza ortonormalna jej wektorów
własnych i rzeczywiste liczby $\lambda_1(\xi),\ldots,\lambda_m(\xi)$.
W tej bazie $\det(\tau I_m+A(\xi))=\prod_k(\tau+\lambda_k(\xi))$.
Dla konormalnej do płaszczyzny $t=\mathrm{const}$ mamy $(\tau,\xi)=(1,0)$,
a zatem $P=I_m$ jest odwracalna.
\end{proof}

\begin{uwaga}
Symetryczna hiperboliczność nie wymaga, aby wartości
własne macierzy $A(\xi)=\sum_j\xi_jA_j$ były różne.
Ścisła hiperboliczność względem czasu oznacza, że
dla każdego $\xi\ne0$ wielomian charakterystyczny
w zmiennej $\tau$ ma $m$ różnych pierwiastków
rzeczywistych (w każdym rozpatrywanym punkcie).
Nie należy utożsamiać tych dwóch własności. Co więcej,
sama rzeczywistość pierwiastków wielomianu charakterystycznego nie
zapewnia oszacowania energetycznego; zobaczymy to w
kontrprzykładzie~\ref{ex:pde21-jordan}.
\end{uwaga}

\section{Tożsamość energii i zależność od danych}
\begin{lemat}[Lokalny bilans energii]\label{lem:pde21-local}
Niech $A_j\in C^1([0,T]\times\R^d;\R^{m\times m})$ będą symetryczne,
$B$ ciągła, a $u\in C^1([0,T]\times\R^d;\R^m)$ rozwiązuje
\eqref{eq:pde21-main}. Wówczas punktowo zachodzi
\begin{equation}\label{eq:pde21-flux}
 \partial_t|u|^2+\sum_{j=1}^d\partial_{x_j}(u^\mathsf TA_ju)
 =u^\mathsf T\!\left(\sum_{j=1}^d\partial_{x_j}A_j-B-B^\mathsf T\right)u
 +2u^\mathsf Tf.
\end{equation}
\end{lemat}
\begin{proof}
Mnożymy równanie przez $2u^\mathsf T$. Otrzymujemy
$2u^\mathsf Tu_t+2\sum_j u^\mathsf TA_ju_{x_j}+2u^\mathsf TBu=2u^\mathsf Tf$.
Pierwszy składnik jest $\partial_t|u|^2$. Symetria $A_j$ daje
$\partial_{x_j}(u^\mathsf TA_ju)=2u^\mathsf TA_ju_{x_j}
+u^\mathsf T(\partial_{x_j}A_j)u$.
Wreszcie $2u^\mathsf TBu=u^\mathsf T(B+B^\mathsf T)u$.
Przenosząc odpowiednie wyrazy, otrzymujemy~\eqref{eq:pde21-flux}.
\end{proof}

\begin{twierdzenie}[Oszacowanie energii na całej przestrzeni]\label{thm:pde21-energy}
Niech $A_j$ będą wspólnie mierzalnymi macierzami
symetrycznymi prawie wszędzie, o jednostajnie ograniczonych
normach w $W^{1,\infty}(\R^d;\R^{m\times m})$
dla prawie każdego $t$ (nie wymagamy pochodnej czasowej),
$B\in L^\infty((0,T)\times\R^d;\R^{m\times m})$,
$f\in L^1(0,T;H)$ oraz $u_0\in H$.
Niech $u\in C([0,T];H)\cap L^1(0,T;H^1(\R^d;\R^m))$
rozwiązuje~\eqref{eq:pde21-main} w sensie dystrybucyjnym
i spełnia $u(0)=u_0$ w $H$.
Oznaczmy
\[
 D=\sum_{j=1}^d\partial_{x_j}A_j-B-B^\mathsf T,\qquad
 M=\operatorname*{ess\,sup}_{(t,x)}
       \|D(t,x)\|_{\mathrm{op}}.
\]
Pochodne współczynników rozumiemy w sensie słabym.
Wtedy dla każdego $t\in[0,T]$
\begin{equation}\label{eq:pde21-L2}
 \|u(t)\|_2\le e^{Mt/2}\left(\|u_0\|_2+
 \int_0^t\|f(r)\|_2\,\mathrm d\lambda_1(r)\right).
\end{equation}
To samo oszacowanie zachodzi dla różnicy dwóch rozwiązań, z prawą
stroną równą różnicy wymuszeń.
\end{twierdzenie}
\begin{proof}
Z równania dystrybucyjnego i ograniczoności
współczynników wynika
\[
 u'=f-\sum_j A_j\partial_{x_j}u-Bu\in L^1(0,T;H).
\]
Zatem $u$ należy do $W^{1,1}(0,T;H)$,
a jego absolutnie ciągły reprezentant zgadza się
z przyjętym ciągłym reprezentantem.
Dla $E(t)=\|u(t)\|_2^2$ reguła łańcuchowa w przestrzeni
Hilberta daje $E'(t)=2(u'(t),u(t))_H$ prawie wszędzie.
Dla prawie każdego $t$ mamy $u(t)\in H^1$,
a iloczyn $u^\mathsf TA_ju$ należy do $W^{1,1}(\R^d)$.
Istotnie, aproksymacja $u(t)$ przez funkcje z
$C_c^\infty$ w $H^1$, ograniczoność $A_j$ i jego
pochodnych oraz nierówność Cauchy'ego--Schwarza
dają zbieżność iloczynów w $W^{1,1}$, z pochodną
\[
 \partial_{x_j}(u^\mathsf TA_ju)
 =2u^\mathsf TA_j\partial_{x_j}u+
 u^\mathsf T(\partial_{x_j}A_j)u.
\]
Wybierzmy $\eta\in C_c^\infty(\R^d)$, $0\le\eta\le1$,
równą $1$ w otoczeniu zera, i połóżmy
$\eta_R(x)=\eta(x/R)$. Całkowanie przez części daje
\[
 \int_{\R^d}\eta_R\partial_{x_j}(u^\mathsf TA_ju)
 \,\mathrm d\lambda_d(x)
 =-\int_{\R^d}(\partial_{x_j}\eta_R)u^\mathsf TA_ju
 \,\mathrm d\lambda_d(x).
\]
Moduł prawej strony jest nie większy niż
$R^{-1}\|\nabla\eta\|_\infty\|A_j\|_\infty\|u(t)\|_2^2$,
a lewa zbiega do całki po całej przestrzeni
z twierdzenia o zbieżności zmajoryzowanej.
W konsekwencji całka pochodnej jest zerowa.
Po podstawieniu równania do pochodnej $E$
otrzymujemy globalny odpowiednik bilansu
~\eqref{eq:pde21-flux}:
\[
 E'(t)=\int_{\R^d}u(t,x)^\mathsf TD(t,x)u(t,x)
 \,\mathrm d\lambda_d(x)+2(u(t),f(t))_H.
\]
Wszystkie składniki są całkowalne w czasie,
ponieważ $u$ jest ograniczona w $H$, $D$ jest
ograniczona, a $f\in L^1(0,T;H)$.
Stąd $E'\le ME+2\sqrt E\|f\|_2$ prawie wszędzie.
Aby nie dzielić przez zero, definiujemy
$y_\eps(t)=\sqrt{E(t)+\eps^2}$ dla $\eps>0$.
Z absolutnej ciągłości i reguły łańcuchowej mamy
\[
 y_\eps'=\frac{E'}{2y_\eps}
 \le\frac{M}{2}y_\eps+\|f\|_2.
\]
Mnożymy przez $e^{-Mt/2}$ i całkujemy od $0$ do $t$:
\[
 y_\eps(t)\le e^{Mt/2}y_\eps(0)
 +\int_0^t e^{M(t-r)/2}\|f(r)\|_2\,\mathrm d\lambda_1(r).
\]
Przechodzimy z $\eps\downarrow0$. Ponieważ $M\ge0$,
dostajemy~\eqref{eq:pde21-L2}. Dla różnicy dwóch
rozwiązań z tej samej klasy stosujemy ten sam dowód
z różnicą danych i wymuszeń.
\end{proof}

\begin{wniosek}[Jednoznaczność]
W klasie rozwiązań z twierdzenia~\ref{thm:pde21-energy}
zagadnienie Cauchy'ego dla~\eqref{eq:pde21-main}
ma co najwyżej jedno rozwiązanie.
\end{wniosek}
\begin{proof}
Dla różnicy dwóch rozwiązań z tymi samymi danymi mamy zerowe
dane początkowe i zerową prawą stronę. Stosujemy~\eqref{eq:pde21-L2}.
\end{proof}

\section{Stożki zależności i skończona prędkość propagacji}
\begin{definicja}[Maksymalna prędkość symbolu]
Dla symetrycznych $A_j$ definiujemy
$c_*:=\operatorname*{ess\,sup}_{(t,x)\in(0,T)\times\R^d}
\sup_{|\nu|=1}\|\sum_j\nu_jA_j(t,x)\|_{\mathrm{op}}$.
Supremum po $\nu$ jest mierzalną funkcją $(t,x)$,
gdyż ciągłość względem $\nu$ pozwala zastąpić sferę
jednostkową jej przeliczalnym gęstym podzbiorem.
Dla symetrycznej macierzy $A(\nu)=\sum_j\nu_jA_j$
norma operatorowa jest maksimum modułów jej wartości
własnych. W szczególności
$c_*\le\sum_j\|A_j\|_{L^\infty;\mathrm{op}}$.
W dalszych rozważaniach zakładamy $c_*<\infty$.
Jeżeli współczynniki są ciągłe na $[0,T]\times\R^d$,
to nierówność $\|A(\nu)\|_{\mathrm{op}}\le c_*$
zachodzi w każdym punkcie i dla każdego $|\nu|=1$:
przekroczenie tej granicy dla ustalonego $\nu$
utrzymywałoby się na względnie otwartym zbiorze
dodatniej miary w walcu.
\end{definicja}

\begin{twierdzenie}[Stożkowa zasada zależności]\label{thm:pde21-cone}
Niech $A_j\in C^1_b([0,T]\times\R^d)$ będą symetryczne,
$B\in C_b([0,T]\times\R^d)$ i niech $u$ będzie klasycznym
rozwiązaniem~\eqref{eq:pde21-main}. Ustalmy $x_0\in\R^d$,
$t_0\in(0,T]$, $R\ge0$ i $c\ge c_*$. Jeśli
$u(0,x)=0$ dla $|x-x_0|\le R+ct_0$ oraz
$f(t,x)=0$ dla $0\le t\le t_0$, $|x-x_0|\le R+c(t_0-t)$,
to $u(t_0,x)=0$ dla $|x-x_0|\le R$.
Dla $R=0$ wniosek dotyczy punktu $x_0$.
\end{twierdzenie}
\begin{proof}
Najpierw załóżmy $R>0$. Niech $r(t)=R+c(t_0-t)$
i $D(t)=B_{r(t)}(x_0)$. Funkcję
$E_D(t)=\int_{D(t)}|u(t,x)|^2\,\mathrm d\lambda_d(x)$
zapisujemy na ustalonej kuli przez podstawienie
$x=x_0+r(t)y$:
\[
 E_D(t)=r(t)^d\int_{B_1(0)}
 g(t,x_0+r(t)y)\,\mathrm d\lambda_d(y),\qquad
 g(t,x)=|u(t,x)|^2.
\]
Promień jest dodatni, a $g$ i jej pierwsze pochodne
są ciągłe i ograniczone na rozpatrywanym zwartym
walcu. Różniczkowanie pod znakiem całki i powrót
do zmiennej $x$ dają
\[
 \begin{split}
 E_D'(t)
 &=\int_{D(t)}\partial_tg(t,x)\,\mathrm d\lambda_d(x)\\
 &\quad+\frac{r'(t)}{r(t)}
 \int_{D(t)}\operatorname{div}_x((x-x_0)g(t,x))
 \,\mathrm d\lambda_d(x).
 \end{split}
\]
Na brzegu kuli $(x-x_0)\cdot\nu=r(t)$, gdzie
$\nu$ jest normalną jednostkową zewnętrzną. Z twierdzenia
Gaussa otrzymujemy więc wzór transportu Reynoldsa
w tym szczególnym przypadku:
\[
 E_D'(t)=\int_{D(t)}\partial_tg(t,x)\,\mathrm d\lambda_d(x)
 +r'(t)\int_{\partial D(t)}g(t,x)\,\mathrm d\mathcal H^{d-1}(x).
\]
Dla $d=1$ to samo obliczenie jest regułą różniczkowania
całki o ruchomych końcach, a $\mathcal H^0$ jest miarą
liczącą na dwóch końcach przedziału.

Oznaczmy $D_A=\sum_j\partial_{x_j}A_j-B-B^\mathsf T$
i $M=\|D_A\|_{L^\infty;\mathrm{op}}$.
Wstawiając lokalny bilans~\eqref{eq:pde21-flux}
oraz $r'=-c$ i stosując twierdzenie Gaussa do strumienia,
otrzymujemy
\[
 \begin{split}
 E_D'(t)
 &=\int_{D(t)}(u^\mathsf TD_Au+2u^\mathsf Tf)
       \,\mathrm d\lambda_d(x)\\
 &\quad-\int_{\partial D(t)}
       (u^\mathsf TA_\nu u+c|u|^2)\,\mathrm d\mathcal H^{d-1}(x),
 \end{split}
\]
gdzie $A_\nu=\sum_j\nu_jA_j$. Z punktowego ograniczenia
$\|A_\nu\|_{\mathrm{op}}\le c_*\le c$ wynika
$u^\mathsf TA_\nu u+c|u|^2\ge0$. Całka brzegowa
w bilansie ma więc znak niedodatni. Ponieważ
$f=0$ w rozważanym stożku, otrzymujemy
$E_D'(t)\le M E_D(t)$. Dane początkowe dają $E_D(0)=0$.
Funkcja $e^{-Mt}E_D(t)$ jest nieujemna i nierosnąca,
zatem $E_D(t)=0$ dla $0\le t\le t_0$.
W otwartej kuli $D(t)$ funkcja $u(t,\cdot)$ znika
punktowo: niezerowa wartość przez ciągłość dawałaby
otoczenie dodatniej miary z dodatnią energią.
Przez ciągłość znika również na brzegu kuli. Dla $R=0$ i $c>0$ przeprowadzamy ten sam rachunek na $[0,t_0)$:
promień $r(t)=c(t_0-t)$ jest tam dodatni, a zerowa energia daje
$u(t,x_0)=0$ dla wszystkich $t<t_0$. Ciągłość w czasie pozwala
przejść do $t=t_0$. Jeżeli $c=0$, to również $c_*=0$, a więc
$A_j=0$ i równanie w każdym punkcie jest zwykłym równaniem
różniczkowym. Istotnie, podstawienie $\nu=e_j$
daje $A_j=0$ w każdym punkcie. Na prostej $x=x_0$
mamy $u_t=-Bu$, gdyż $f(t,x_0)=0$, i $u(0,x_0)=0$.
Nierówność $(|u|^2)'\le2\|B\|_\infty |u|^2$
daje zerowanie także w tym przypadku.
\end{proof}

\begin{wniosek}[Lokalna zależność od danych i propagacja nośnika]
Dwa klasyczne rozwiązania tego samego układu
(z różnymi danymi początkowymi i wymuszeniami)
mają równe wartości na $\{x:|x-x_0|\le R\}$ w chwili
$t_0$, jeśli ich dane początkowe i wymuszenia są równe
w zbiorach wskazanych w twierdzeniu~\ref{thm:pde21-cone}.
W szczególności przy $f=0$ i zwartym nośniku danych
początkowych $K$ zachodzi
\[
 \operatorname{supp}u(t,\cdot)
 \subset \{x\in\R^d:\operatorname{dist}(x,K)\le ct\},
 \qquad 0\le t\le T,\quad c\ge c_*.
\]
\end{wniosek}
\begin{proof}
Pierwszą część otrzymujemy, stosując twierdzenie
~\ref{thm:pde21-cone} do różnicy rozwiązań.
Dla $t>0$ i punktu $x$ o $\operatorname{dist}(x,K)>ct$
zbiór $\{y:|y-x|\le ct\}$ nie przecina $K$,
więc twierdzenie z $R=0$ daje $u(t,x)=0$.
Zbiór takich punktów jest otwarty (funkcja odległości
od niepustego zbioru jest lipschitzowska ze stałą $1$
na mocy nierówności trójkąta), więc nie przecina nośnika.
Dla $t=0$ wniosek wynika z definicji $K$.
Jeżeli dane są identycznie zerowe, twierdzenie zastosowane
w każdym punkcie daje $u=0$; przyjmujemy wtedy nośnik
pusty i odległość od niego równą $+\infty$.
\end{proof}



\begin{center}
\begin{tikzpicture}[scale=.88,>=Latex]
\fill[blue!8] (0.8,0) -- (6.0,0) -- (3.4,3.55) -- cycle;
\draw[->] (0,0) -- (6.8,0) node[right] {$x$};
\draw[->] (0,0) -- (0,4.2) node[above] {$t$};
\draw[thick,blue!65!black] (.8,0)--(3.4,3.55)--(6,0);
\draw[dashed] (1.6789,1.2)--(5.1211,1.2);
\draw[gray,dotted] (0,3.55)--(3.4,3.55);
\draw[gray,dotted] (3.4,0)--(3.4,3.55);
\draw[<->,blue!65!black] (.8,-.25)--(3.4,-.25)
 node[midway,below] {$ct_0$};
\draw[<->,blue!65!black] (3.4,-.25)--(6,-.25)
 node[midway,below] {$ct_0$};
\node[left] at (0,3.55) {$t_0$};
\node[below] at (3.4,-.65) {$x_0$};
\node[above right] at (3.4,3.55) {$(t_0,x_0)$};
\node[align=center,font=\small] at (3.4,.55) {dane i wymuszenie\\w stożku};
\node[right] at (5.1,1.2) {$t=t_1$};
\end{tikzpicture}

\small Rysunek: wsteczny stożek zależności w jednym wymiarze.
Gdy $R=0$, podstawa ma promień $ct_0$.
\end{center}

\section{Równanie falowe jako układ symetryczny}
Niech $c>0$ i niech $u\in C^2([0,T]\times\R)$
spełnia $u_{tt}-c^2u_{xx}=0$. Określamy $v=u_t$,
$w=cu_x$ i $U=(v,w)^\mathsf T$.
Z równania falowego mamy $v_t=cw_x$, a z równości
pochodnych mieszanych $w_t=cv_x$. Wówczas
\begin{equation}\label{eq:pde21-wavefirst}
 U_t+\begin{pmatrix}0&-c\\-c&0\end{pmatrix}U_x=0,
 \qquad |U|^2=u_t^2+c^2u_x^2.
\end{equation}
Macierz główna jest symetryczna; wektory
$2^{-1/2}(1,-1)^\mathsf T$ i $2^{-1/2}(1,1)^\mathsf T$
są jej wektorami własnymi odpowiednio dla $c$ i $-c$.
Są to prędkości przenoszenia dwóch niezależnych modów.

Odwrotnie, niech $(v,w)\in C^1([0,T]\times\R;\R^2)$
spełnia~\eqref{eq:pde21-wavefirst}.
Jeżeli $a\in C^2(\R)$ i $w(0,x)=ca'(x)$, to definiujemy
\[
 u(t,x)=a(x)+\int_0^t v(r,x)\,\mathrm d\lambda_1(r).
\]
Z twierdzenia Newtona--Leibniza mamy $u_t=v$.
Różniczkując pod znakiem całki względem $x$ i używając
$w_t=cv_x$, otrzymujemy
\[
 cu_x(t,x)=ca'(x)+\int_0^t w_t(r,x)\,\mathrm d\lambda_1(r)
 =w(t,x).
\]
Obie pierwsze pochodne $u$ są klasy $C^1$, więc
$u\in C^2$ i $u_{tt}=v_t=cw_x=c^2u_{xx}$.
Dane początkowe to $u(0,x)=a(x)$ i $u_t(0,x)=v(0,x)$.
Wybór $a$ jest jednoznaczny z dokładnością do
stałej addytywnej, jeśli znamy tylko $w(0,\cdot)$:
można wziąć $a(x)=a(0)+c^{-1}\int_0^x w(0,y)\,
\mathrm d\lambda_1(y)$. Układ opisuje więc pochodne
rozwiązania falowego, ale sam nie wyznacza jego
stałej addytywnej.

\begin{stwierdzenie}[Rozkład na fale podróżujące]
Niech $v,w$ będą gładkimi funkcjami rozwiązującymi
\eqref{eq:pde21-wavefirst} na $[0,T]\times\R$. Funkcje
$r_+=v-w$ i $r_-=v+w$ spełniają odpowiednio
$(\partial_t+c\partial_x)r_+=0$ oraz
$(\partial_t-c\partial_x)r_-=0$.
W konsekwencji $r_+(t,x)=r_+(0,x-ct)$ i
$r_-(t,x)=r_-(0,x+ct)$.
\end{stwierdzenie}
\begin{proof}
Równania składowych to $v_t=cw_x$ i $w_t=cv_x$.
Po odjęciu otrzymujemy $(v-w)_t=-c(v-w)_x$,
a po dodaniu $(v+w)_t=c(v+w)_x$.
Dla ustalonego $(t,x)$ funkcja
$r_+(s,x-c(t-s))$ ma pochodną
$(\partial_t+c\partial_x)r_+$ w punkcie
$(s,x-c(t-s))$, czyli zero. Porównując jej wartości
dla $s=0$ i $s=t$, otrzymujemy pierwszą formułę.
Analogicznie funkcja $r_-(s,x+c(t-s))$ ma pochodną
$(\partial_t-c\partial_x)r_-=0$, co daje drugą.
Jeżeli $u(0,x)=a(x)$ i $u_t(0,x)=b(x)$, to
\[
 r_+(0,x)=b(x)-ca'(x),\qquad
 r_-(0,x)=b(x)+ca'(x).
\]
Z $v=(r_++r_-)/2$ i $w=(r_--r_+)/2$
odzyskujemy obie pochodne rozwiązania falowego.
Całkując $v$ względem czasu z danymi $a$, dostajemy
\[
 u(t,x)=\frac{a(x-ct)+a(x+ct)}{2}
 +\frac1{2c}\int_{x-ct}^{x+ct}b(y)\,\mathrm d\lambda_1(y),
\]
gdyż podstawienia $y=x\mp cr$ w całkach z $b(x\mp cr)$
dają drugi składnik, a całki z $a'$ dają pierwszy
po dodaniu $a(x)$. Jest to formuła d'Alemberta
z rozdziału~\ref{chap:pde-wave}.
\end{proof}

\section{Oszacowania wyższych pochodnych bez ich utraty}
\begin{lemat}[Komutator z pochodną wielokrotną]\label{lem:pde21-comm}
Niech $s\ge1$ będzie liczbą całkowitą,
$A\in W^{s,\infty}(\R^d;\R^{m\times m})$ oraz
$u\in H^s(\R^d;\R^m)$. Dla każdego wielowskaźnika
$\alpha$ z $|\alpha|\le s$ zachodzi
\[
 \|[\partial^\alpha,A]\partial_j u\|_2
 \le C_{d,s}\|A\|_{W^{s,\infty}}\|u\|_{H^s},
\]
gdzie komutator oznacza
$[\partial^\alpha,A]v=\partial^\alpha(Av)-A\partial^\alpha v$.
Dla $v=\partial_j u$ różnicę rozumiemy początkowo
w $H^{-1}(\R^d;\R^m)$: mnożenie przez $A$ na $H^{-1}$
definiujemy przez dualność, używając ograniczonego
operatora mnożenia przez $A^\mathsf T$ na $H^1$.
Teza stwierdza, że różnica należy faktycznie do $H$.
\end{lemat}
\begin{proof}
Najpierw dla $u\in C_c^\infty$ stosujemy słabą
regułę Leibniza (przez iterację reguły iloczynu
funkcji Sobolewa z funkcją gładką):
\[
 [\partial^\alpha,A]\partial_ju
 =\sum_{0<\beta\le\alpha}\binom{\alpha}{\beta}
 (\partial^\beta A)\partial_j\partial^{\alpha-\beta}u,
 \qquad
 \binom{\alpha}{\beta}=\prod_{k=1}^d
 \binom{\alpha_k}{\beta_k}.
\]
Nierówność $\beta\le\alpha$ rozumiemy współrzędnie.
Dla $\alpha=0$ suma jest pusta. Dla każdego
pozostałego składnika mamy
$1+|\alpha|-|\beta|\le|\alpha|\le s$.
Norma $L^2$ tego składnika nie przekracza
$\binom{\alpha}{\beta}\|\partial^\beta A\|_\infty
\|\partial_j\partial^{\alpha-\beta}u\|_2$.
Sumowanie daje żądane oszacowanie.
Dla ogólnego $u\in H^s$ wybieramy $u_k\in C_c^\infty$
zbieżne w $H^s$. Prawa strona wzoru zbiega w $H$
na mocy powyższych oszacowań. Jednocześnie
$A\partial_j u_k\to A\partial_j u$ w $H^{s-1}$:
słaba reguła Leibniza i ograniczoność pochodnych $A$
dają ciągłość mnożenia na $H^{s-1}$.
Zatem pierwszy składnik komutatora zbiega w $H^{-1}$.
Drugi zbiega w $H^{-1}$ dzięki ciągłości mnożenia
na tej przestrzeni. Ponieważ zanurzenie $H\hookrightarrow H^{-1}$
jest ciągłe i iniektywne, granica w $H$ jest właśnie
żądanym komutatorem. Nie przybliżamy $A$ w normie
$W^{s,\infty}$.
\end{proof}

\begin{twierdzenie}[Energia w $H^s$]\label{thm:pde21-Hs}
Niech $s\ge1$ będzie liczbą całkowitą,
$A_j,B\in C([0,T];W^{s,\infty}(\R^d;\R^{m\times m}))$,
a macierze $A_j$ niech będą symetryczne.
Przyjmijmy normę
$\|u\|_{H^s}^2=\sum_{|\alpha|\le s}\|\partial^\alpha u\|_2^2$.
Niech $u_0\in H^s$, $f\in L^1(0,T;H^s)$
i niech $u\in C([0,T];H^s)\cap L^1(0,T;H^{s+1})$
będzie rozwiązaniem~\eqref{eq:pde21-main} w sensie
dystrybucyjnym, z danymi początkowymi $u_0$.
Wówczas
\begin{equation}\label{eq:pde21-Hs}
 \|u(t)\|_{H^s}\le e^{Ct}
 \left(\|u_0\|_{H^s}+\int_0^t\|f(r)\|_{H^s}\,\mathrm d\lambda_1(r)\right),
\end{equation}
gdzie $C$ zależy od $d,m,s$ i jednostajnych norm
$W^{s,\infty}$ współczynników, nie zaś od $u$ i $f$.
\end{twierdzenie}
\begin{proof}
Niech $|\alpha|\le s$. Po zastosowaniu $\partial^\alpha$
do~\eqref{eq:pde21-main} otrzymujemy
\[
 \partial_t(\partial^\alpha u)+\sum_jA_j\partial_j(\partial^\alpha u)
 +B\partial^\alpha u
 =\partial^\alpha f-\sum_j[\partial^\alpha,A_j]\partial_ju
 -[\partial^\alpha,B]u.
\]
Oznaczmy $v_\alpha=\partial^\alpha u$ i
$R_\alpha=-\sum_j[\partial^\alpha,A_j]\partial_ju
-[\partial^\alpha,B]u$.
Z lematu~\ref{lem:pde21-comm} oraz słabej reguły
Leibniza dla $B$ (tu każdy składnik zawiera co
najwyżej $s-1$ pochodnych $u$) wynika
\[
 \left(\sum_{|\alpha|\le s}\|R_\alpha(t)\|_2^2\right)^{1/2}
 \le C_0\|u(t)\|_{H^s},
\]
gdzie $C_0$ zależy wyłącznie od podanych norm
współczynników i liczby wielowskaźników.
Każde $v_\alpha$ należy do $C([0,T];H)\cap L^1(0,T;H^1)$,
a wymuszenie $\partial^\alpha f+R_\alpha$ do
$L^1(0,T;H)$. Stosujemy więc globalną tożsamość
energii z dowodu twierdzenia~\ref{thm:pde21-energy}.
Po zsumowaniu i zastosowaniu nierówności
Cauchy'ego--Schwarza do skończonej sumy iloczynów
skalarnych otrzymujemy dla
$E_s=\sum_{|\alpha|\le s}\|v_\alpha\|_2^2$
\[
 E_s'\le M E_s+2\sqrt{E_s}
   \bigl(\|f\|_{H^s}+C_0\sqrt{E_s}\bigr).
\]
Zatem $E_s'\le(M+2C_0)E_s+2\sqrt{E_s}\|f\|_{H^s}$.
Argument z $\sqrt{E_s+\eps^2}$ i czynnikiem całkującym
z dowodu twierdzenia~\ref{thm:pde21-energy}
daje~\eqref{eq:pde21-Hs} z $C=M/2+C_0$.
Poniższy lemat pozwoli usunąć pomocniczy warunek
całkowalności w $H^{s+1}$.
\end{proof}


\begin{lemat}[Przedłużenie tożsamości energii na rozwiązania słabe]
\label{lem:pde21-friedrichs-comm}
Niech $A_j,B$ spełniają założenia dotyczące
współczynników z twierdzenia~\ref{thm:pde21-energy}.
Jeżeli $v\in L^\infty(0,T;H)$ rozwiązuje dystrybucyjnie
$v_t+\sum_j A_j\partial_jv+Bv=F$ z $F\in L^1(0,T;H)$,
to ma jednoznacznego reprezentanta w $C([0,T];H)$.
Dla każdego $0\le a\le b\le T$ zachodzi
\[
 \begin{split}
 \|v(b)\|_2^2-\|v(a)\|_2^2
 &=\int_a^b\int_{\R^d}v(t,x)^\mathsf TD_A(t,x)v(t,x)
       \,\mathrm d\lambda_d(x)\,\mathrm d\lambda_1(t)\\
 &\quad+2\int_a^b(v(t),F(t))_H\,\mathrm d\lambda_1(t),
 \end{split}
\]
gdzie $D_A=\sum_j\partial_{x_j}A_j-B-B^\mathsf T$.
W szczególności zachodzi oszacowanie~\eqref{eq:pde21-L2}
z dowolnej chwili początkowej.
\end{lemat}
\begin{proof}
\emph{Krok 1: słabo ciągły reprezentant.}
Iloczyn z pochodną dystrybucyjną określamy przez
$A_j\partial_j v=\partial_j(A_jv)-(\partial_jA_j)v$,
co należy do $H^{-1}$ dla prawie każdego $t$.
Z równania wynika
$v'\in L^1(0,T;H^{-1})$, ponieważ $\partial_j:H\to H^{-1}$
jest ograniczony, współczynniki mają jednostajne
ograniczenia, a $F\in L^1(0,T;H)$.
Wybieramy absolutnie ciągły reprezentant
$w:[0,T]\to H^{-1}$ klasy $v\in W^{1,1}(0,T;H^{-1})$.
Niech $K=\|v\|_{L^\infty(0,T;H)}$ i niech $G$
będzie zbiorem czasów pełnej miary, dla których
$w(t)=v(t)$ w $H^{-1}$ oraz $\|v(t)\|_2\le K$.
Dla każdego $t$ wybieramy $t_n\in G$ z $t_n\to t$.
Ciąg $v(t_n)$ ma słabo zbieżny w $H$ podciąg
z granicą $h_t$, $\|h_t\|_2\le K$. Z ciągłości
zanurzenia $H\hookrightarrow H^{-1}$ jego słaba
granica w $H^{-1}$ jest obrazem $h_t$, ale
$w(t_n)\to w(t)$ silnie w $H^{-1}$.
Stąd $w(t)=h_t$ po utożsamieniu przez zanurzenie.
Jego iniektywność zapewnia jednoznaczność $h_t$.
Przyjmijmy $v(t)=h_t$ dla wszystkich $t$.
Dla $\phi\in H^1$ funkcja $t\mapsto(v(t),\phi)_H$
jest ciągła, bo jest wartością dualności $H^{-1},H^1$.
Aproksymacja dowolnego $\phi\in H$ funkcjami z $H^1$
i jednolite ograniczenie $\|v(t)\|_2\le K$ dają
słabą ciągłość $v:[0,T]\to H$.

\emph{Krok 2: komutator wygładzania.}
Wybierzmy $\rho\in C_c^\infty(\R^d)$ nieujemną,
o całce $1$, i połóżmy
$\rho_\eps(z)=\eps^{-d}\rho(z/\eps)$,
$J_\eps v=\rho_\eps*v$, gdzie splot jest tylko
względem zmiennych przestrzennych.
Operator $J_\eps$ jest kontrakcją w $H$ oraz
$J_\eps g\to g$ w $H$ dla każdego $g\in H$.
Funkcja $v_\eps=J_\eps v$ spełnia
\[
 \partial_t v_\eps+\sum_j A_j\partial_j v_\eps+Bv_\eps
 =J_\eps F+R_\eps,\qquad
 R_\eps=\sum_j[A_j,J_\eps]\partial_jv+[B,J_\eps]v,
\]
gdzie $[A,J_\eps]=AJ_\eps-J_\eps A$.
Dla ustalonego $t$ i początkowo gładkiej $v$
\[
 [J_\eps,A]\partial_jv(x)=
 \int_{\R^d}\rho_\eps(z)(A(x-z)-A(x))\partial_jv(x-z)
 \,\mathrm d\lambda_d(z).
\]
Całkując przez części względem $z_j$ (pochodna
$v(x-z)$ względem $z_j$ ma przeciwny znak),
otrzymujemy
\[
 \int_{\R^d}
 \bigl(\partial_j\rho_\eps(z)(A(x-z)-A(x))
 -\rho_\eps(z)(\partial_j A)(x-z)\bigr)v(x-z)
 \,\mathrm d\lambda_d(z).
\]
Całka samego współczynnika przy $v(x-z)$ jest
zerowa: jest całką pochodnej
$\partial_{z_j}(\rho_\eps(z)(A(x-z)-A(x)))$.
Możemy więc zastąpić $v(x-z)$ przez $v(x-z)-v(x)$.
Dla $A\in W^{1,\infty}$ używamy lipschitzowskiego
reprezentanta; powyższe całkowanie przez części
jest dopuszczalne w sensie słabym.
Przy $\operatorname{supp}\rho\subset\overline{B_{R_\rho}(0)}$
nierówność Minkowskiego dla całki daje
\[
 \|[J_\eps,A]\partial_jv\|_2
 \le C_\rho\|\nabla A\|_\infty
       \sup_{|z|\le\eps R_\rho}\|v(\cdot-z)-v\|_2,
\]
gdzie $C_\rho$ można wybrać jako
$\int |z|\,|\partial_j\rho(z)|\,\mathrm d\lambda_d(z)
+\|\rho\|_{L^1}$, z odpowiednią normą pochodnych
macierzy. Stała nie zależy od $\eps$, gdyż podstawienie
$z=\eps y$ znosi czynniki skalowania w pierwszej całce.
Wzór i oszacowanie rozszerzają się na $v\in H$ przez
aproksymację w $H$: prawa strona określa ograniczony
operator na $H$, a lewa zbiega dystrybucyjnie dzięki
definicji $A\partial_j v$ i ciągłości splotu.
Ciągłość translacji w $H$ daje zbieżność do zera
dla prawie każdego $t$, a oszacowanie przez
$2C_\rho\|\nabla A\|_\infty K$ i twierdzenie
o zbieżności zmajoryzowanej dają zbieżność
w $L^1(0,T;H)$. Dla członu bez pochodnej piszemy
\[
 [B,J_\eps]v=B(J_\eps v-v)+(Bv-J_\eps(Bv)).
\]
Oba składniki zbiegają do zera w $L^1(0,T;H)$
z silnej zbieżności splotu, ograniczoności $B$
i tego samego twierdzenia granicznego.
W konsekwencji $\|R_\eps\|_{L^1(0,T;H)}\to0$.
Również $J_\eps F\to F$ w $L^1(0,T;H)$.

\emph{Krok 3: bilans dla wszystkich chwil.}
Dla ustalonego $\eps$ operator $J_\eps$ jest
ograniczony z $H^{-1}$ do $H^k$ dla każdego
całkowitego $k\ge0$ (wynika to np. z szybkiego zaniku
transformaty Fouriera $\rho_\eps$).
Stąd $v_\eps=J_\eps w$ jest ciągły w $H^1$
na całym $[0,T]$ i należy do $W^{1,1}(0,T;H)$.
Jego wartości w każdej chwili są dokładnie
$J_\eps v(t)$, a nie tylko prawie wszędzie.
Stosujemy tożsamość z dowodu twierdzenia
~\ref{thm:pde21-energy} do równania wygładzonego,
z wymuszeniem $J_\eps F+R_\eps$, i całkujemy
od dowolnego $a$ do dowolnego $b$.
Dla każdego ustalonego $t\in[0,T]$ mamy
$J_\eps v(t)\to v(t)$ silnie w $H$ na mocy
własności aproksymacji jedynki, także w chwilach
spoza zbioru $G$. W szczególności obie wartości
końcowe w bilansie zbiegają do właściwych norm.
Ponadto $\|v_\eps(t)\|_2\le K$, a różnica wyrazów
kwadratowych z $D_A$ jest ograniczona przez
$2MK\|v_\eps(t)-v(t)\|_2$.
Twierdzenie o zbieżności zmajoryzowanej daje
zbieżność tych całek. Dla członu wymuszenia
rozkładamy różnicę na składniki z
$J_\eps F-F+R_\eps$ oraz $v_\eps-v$.
Pierwsze całki mają moduł nie większy niż
$K\|J_\eps F-F+R_\eps\|_{L^1(0,T;H)}$,
druga zbiega do zera przez majorantę $2K\|F(t)\|_2$.
Otrzymujemy zatem tożsamość z tezy dla \emph{każdej}
pary $a,b$, nie tylko prawie wszystkich końców przedziału.

W szczególności funkcja $t\mapsto\|v(t)\|_2^2$
jest absolutnie ciągła, gdyż prawa strona tożsamości
dla $a=0$ jest całką funkcji z $L^1(0,T)$.
Dla $t_n\to t$ słaba ciągłość i ciągłość normy dają
\[
 \|v(t_n)-v(t)\|_2^2
 =\|v(t_n)\|_2^2+\|v(t)\|_2^2-2(v(t_n),v(t))_H\to0.
\]
To dowodzi silnej ciągłości; jednoznaczność
reprezentanta wynika z gęstości zbioru pełnej miary.
Oszacowanie~\eqref{eq:pde21-L2} wynika z tej samej
nierówności różniczkowej dla normy, co w twierdzeniu
~\ref{thm:pde21-energy}, z początkiem całkowania w $a$.
\end{proof}

\section{Istnienie przez regularizację Friedricha}
Warto odróżnić metodę zastosowaną poniżej od rozdziału~20:
w równaniach parabolicznych otrzymywaliśmy zwartość z całki po czasie
normy $H^1$. Tutaj podstawowym narzędziem jest kontrola \emph{wszystkich}
pochodnych przestrzennych dzięki symetrii macierzy głównych.

\begin{twierdzenie}[Zagadnienie Cauchy'ego dla gładkich współczynników]
\label{thm:pde21-existence}
Niech $s>d/2+2$ będzie liczbą całkowitą. Załóżmy, że
$A_j=A_j^\mathsf T$ oraz
$A_j,B\in C([0,T];W^{s+1,\infty}(\R^d;\R^{m\times m}))$,
zaś $u_0\in H^s(\R^d;\R^m)$ i
$f\in C([0,T];H^s(\R^d;\R^m))$.
Wtedy istnieje dokładnie jedno rozwiązanie zagadnienia
\eqref{eq:pde21-main}, dla którego
\[
 u\in C([0,T];H^s(\R^d;\R^m))\cap
 C^1([0,T];H^{s-1}(\R^d;\R^m)).
\]
Spełnia ono~\eqref{eq:pde21-Hs}.
\end{twierdzenie}
\begin{proof}
\emph{Krok 1: regularizacja.}
Niech $P_N$ będzie ortogonalnym rzutem Fouriera na częstotliwości
$|\xi|\le N$. Operator $P_N$ jest kontrakcją w każdym $H^k$
i komutuje z pochodnymi. Szukamy rozwiązania
\begin{equation}\label{eq:pde21-friedrichs}
 u_N'+P_N\left(\sum_j A_j\partial_j u_N+Bu_N\right)=P_Nf,
 \quad u_N(0)=P_Nu_0,
 \quad P_Nu_N=u_N.
\end{equation}
Na podprzestrzeni $P_NH^s$ operator
$v\mapsto P_N(A_j\partial_jv+Bv)$ jest ograniczony, ponieważ
$\|\partial_jv\|_{H^s}\le N\|v\|_{H^s}$.
Równanie~\eqref{eq:pde21-friedrichs} jest więc zwykłym
liniowym równaniem różniczkowym w przestrzeni Banacha, o ciągłych
w czasie współczynnikach; metoda kolejnych przybliżeń daje
rozwiązanie na krótkim przedziale, a liniowość pozwala je
przedłużyć do $T$.

\emph{Krok 2: jednostajne oszacowania.}
Ponieważ $P_N$ jest samosprzężony i $P_N\partial^\alpha u_N=\partial^\alpha u_N$,
rachunek energii po zastosowaniu $\partial^\alpha$ wygląda tak
samo jak w dowodzie~\eqref{eq:pde21-Hs}. Istotnie
$\langle\partial^\alpha u_N,P_N(A_j\partial_j\partial^\alpha u_N)\rangle
=\langle\partial^\alpha u_N,A_j\partial_j\partial^\alpha u_N\rangle$.
Komutatory są oszacowane przez Lemat~\ref{lem:pde21-comm}.
Stąd $\sup_{N,t}\|u_N(t)\|_{H^s}\le C$.
Z~\eqref{eq:pde21-friedrichs} i ograniczoności mnożenia
przez $W^{s,\infty}$ dostajemy także
$\sup_{N,t}\|u_N'(t)\|_{H^{s-1}}\le C$.

\emph{Krok 3: zwartość lokalna.}
Na każdej kuli $B_R$ zanurzenie $H^s(B_R)\hookrightarrow
H^{s-1}(B_R)$ jest zwarte. Ponadto
$\|u_N(t)-u_N(r)\|_{H^{s-1}(B_R)}\le C|t-r|$.
Z twierdzenia Arzelà--Ascoliego w metrycznej przestrzeni
$H^{s-1}(B_R)$ otrzymujemy podciąg zbieżny w
$C([0,T];H^{s-1}(B_R))$. Wybierając podciąg przekątny dla
$R=1,2,\ldots$, dostajemy funkcję $u$, która jest granicą lokalną
oraz słabą-* granicą w $L^\infty(0,T;H^s(\R^d))$.

\emph{Krok 4: równanie i dane.}
Dla $\varphi\in C_c^\infty((0,T)\times\R^d;\R^m)$
mnożymy~\eqref{eq:pde21-friedrichs} przez $\varphi$
i całkujemy. Przenosimy $P_N$ na funkcję testową
$P_N\varphi$; ponieważ $P_N\varphi\to\varphi$ w każdym
$H^k$, przechodzimy do granicy w wyrazach liniowych,
korzystając z lokalnej silnej zbieżności oraz jednostajnych
oszacowań na całej przestrzeni. Otrzymujemy~\eqref{eq:pde21-main}
w sensie dystrybucyjnym. Zbieżność
$u_N(0)=P_Nu_0\to u_0$ w $H^{s-1}$ i lokalna ciągłość implikują
$u(0)=u_0$ jako element $H^{s-1}(\R^d)$.

\emph{Krok 5: energia, jednoznaczność i ciągłość silna.}
Dla każdego wielowskaźnika $|\alpha|\le s$ funkcja
$v=\partial^\alpha u$ należy do $L^\infty(0,T;L^2)$.
Po zróżniczkowaniu równania jej prawa strona jest sumą
$\partial^\alpha f$ i komutatorów z
Lematu~\ref{lem:pde21-comm}, a więc należy do $L^1(0,T;L^2)$.
Lemat~\ref{lem:pde21-friedrichs-comm} stosuje się zatem
osobno do wszystkich $v=\partial^\alpha u$; daje
$v\in C([0,T];L^2)$ i po zsumowaniu oszacowanie
~\eqref{eq:pde21-Hs}. Dzięki temu $u\in C([0,T];H^s)$
bez konieczności zakładania z góry ciągłości najwyższych pochodnych.
Równanie ma prawą stronę ciągłą z wartościami w $H^{s-1}$,
ponieważ mnożenie przez przyjęte współczynniki jest ciągłe
na $H^{s-1}$, zatem $u_t\in C([0,T];H^{s-1})$.
Ostatecznie $u\in C^1([0,T];H^{s-1})$.
Jednoznaczność wynika z Twierdzenia~\ref{thm:pde21-energy}
zastosowanego do różnicy dwóch rozwiązań.
\end{proof}

\begin{uwaga}[Wersja o mniejszej regularności]
Następnie można przez aproksymację danych w $L^2$
i oszacowanie~\eqref{eq:pde21-L2} otrzymać rozwiązania energetyczne
z ciągłością $C([0,T];L^2)$, ale wymaga to osobnego uzasadnienia
reguły energii dla granicznych rozwiązań. Twierdzenie powyżej
celowo sformułowaliśmy z mocniejszymi założeniami i bez ukrytej
utraty pochodnych.
\end{uwaga}

\section{Symetryzatory: gdy macierze nie są symetryczne}
\begin{definicja}[Układ symetryzowalny]\label{def:pde21-symmetrizer}
Niech $A_j$ będą rzeczywistymi macierzami $m\times m$.
Mówimy, że układ~\eqref{eq:pde21-main} ma \emph{symetryzator} $S(t,x)$,
jeżeli $S=S^\mathsf T$, istnieją stałe $0<\sigma\le\Sigma$ takie,
że $\sigma I_m\le S\le\Sigma I_m$, a macierze $SA_j$ są symetryczne.
Dla oszacowań energii zakładamy dodatkowo
$S\in W^{1,\infty}$ i $A_j,B\in W^{1,\infty}$.
\end{definicja}

\begin{twierdzenie}[Energia ważona symetryzatorem]
\label{thm:pde21-symmetrizer}
Niech założenia Definicji~\ref{def:pde21-symmetrizer} zachodzą
na $[0,T]\times\R^d$, a $u$ jest gładkim, zanikającym
w nieskończoności rozwiązaniem~\eqref{eq:pde21-main}.
Wówczas dla $E_S(t)=\int_{\R^d}u^\mathsf TS u\,\mathrm d\lambda_d$
 istnieje stała $C$, zależna tylko od
$\sigma,\Sigma$ i norm odpowiednich współczynników,
taka że
\[
 \sqrt{E_S(t)}\le e^{Ct}\left(\sqrt{E_S(0)}+
 \sqrt\Sigma\int_0^t\|f(r)\|_2\,\mathrm dr\right).
\]
\end{twierdzenie}
\begin{proof}
Różniczkując, otrzymujemy
$\partial_t(u^\mathsf TS u)=2u^\mathsf TS u_t+u^\mathsf TS_tu$.
Po wstawieniu równania i użyciu symetrii $SA_j$:
\[
 \partial_t(u^\mathsf TS u)+\sum_j\partial_j(u^\mathsf TSA_ju)
 =u^\mathsf T\Bigl(S_t+\sum_j\partial_j(SA_j)-SB-B^\mathsf TS\Bigr)u
 +2u^\mathsf TSf.
\]
Całkujemy po przestrzeni. Z $E_S\ge\sigma\|u\|_2^2$
wynika ograniczenie wyrazu kwadratowego przez $C E_S$,
a z nierówności Cauchy'ego--Schwarza w iloczynie skalarnym
wyznaczonym przez $S$ dostajemy
$\left|\int_{\R^d}u(t,x)^\mathsf TS(t,x)f(t,x)\,\mathrm d\lambda_d(x)\right|\le\sqrt{E_S(t)}\sqrt\Sigma\|f(t)\|_2$.
Argument z $\sqrt{E_S+\eps^2}$ i Gronwallem kończy dowód.
\end{proof}

\begin{kontrprzyklad}[Rzeczywiste charakterystyki bez stabilności]\label{ex:pde21-jordan}
Niech $U=(u,v)^\mathsf T$ i rozważmy na
$[0,T]\times\R$ układ $u_t=v_x$, $v_t=0$.
Macierz główna w zapisie $U_t+AU_x=0$ jest równa
$A=\begin{pmatrix}0&-1\\0&0\end{pmatrix}$.
Jej jedyną wartością własną jest $0$, lecz nie ma bazy
wektorów własnych. Wybierzmy niezerową funkcję
$\phi\in C_c^\infty(\R)$ i dane
$u(0,x)=0$, $v(0,x)=n^{-1}\phi(x)\sin(nx)$.
Równanie ma jawne rozwiązanie
\[
 v(t,x)=n^{-1}\phi(x)\sin(nx),\qquad
 u(t,x)=t\bigl(\phi(x)\cos(nx)+n^{-1}\phi'(x)\sin(nx)\bigr).
\]
Norma początkowa $\|(u(0),v(0))\|_{L^2(\R;\R^2)}$
nie przekracza $n^{-1}\|\phi\|_2$ i dąży do zera.
Z kolei dla każdego ustalonego $t>0$,
na mocy lematu Riemanna--Lebesgue'a,
$\|u(t)\|_2^2\to t^2\|\phi\|_2^2/2>0$:
całka $\int_{\R}|\phi(x)|^2\cos^2(nx)\,
\mathrm d\lambda_1(x)$ dąży do $\|\phi\|_2^2/2$,
a składniki z $n^{-1}\phi'$ znikają w $L^2$.
Nie istnieje zatem ograniczony operator ewolucji
$L^2(\R;\R^2)\to L^2(\R;\R^2)$ w żadnej chwili $t>0$.
Przykład na torusie powstaje prościej, biorąc $\phi\equiv1$,
lecz funkcji okresowych nie wolno traktować jako elementów
$L^2(\R)$.
\end{kontrprzyklad}

\section{Zmiennowspółczynnikowe równanie falowe}
Niech $G(x)=G(x)^\mathsf T$ będzie macierzą $d\times d$,
$G\in W^{1,\infty}(\R^d)$ oraz
$\lambda|\xi|^2\le\xi^\mathsf TG(x)\xi\le\Lambda|\xi|^2$.
Rozważamy
\begin{equation}\label{eq:pde21-wavevariable}
 u_{tt}-\operatorname{div}(G(x)\nabla u)=f.
\end{equation}
\begin{twierdzenie}[Energia równania falowego]\label{thm:pde21-wave-energy}
Dla dostatecznie gładkiego rozwiązania~\eqref{eq:pde21-wavevariable},
które wraz z pochodnymi dostatecznie szybko zanika, energia
\[
 \mathcal E(t)=\frac12\int_{\R^d}
 (|u_t|^2+\nabla u^\mathsf TG\nabla u)\,\mathrm d\lambda_d
\]
spełnia $\mathcal E'(t)=\int_{\R^d}f u_t\,\mathrm d\lambda_d$.
Jeżeli $f=0$, energia jest stała. Jeśli $f$ znika we wstecznym
stożku o prędkości $c\ge\sqrt\Lambda$, a dane $u_0,u_1$ znikają
na jego podstawie, to $u$ znika w całym wnętrzu stożka.
\end{twierdzenie}
\begin{proof}
Mnożymy równanie przez $u_t$ i całkujemy.
Ponieważ $G$ jest niezależna od $t$, całkowanie przez części daje
$\int_{\R^d}u_{tt}u_t\,\mathrm d\lambda_d(x)
=\frac12\frac{\mathrm d}{\mathrm dt}\|u_t\|_2^2$
oraz
\[
-\int_{\R^d}\operatorname{div}(G\nabla u)u_t\,\mathrm d\lambda_d(x)
=\int_{\R^d}G\nabla u\cdot\nabla u_t\,\mathrm d\lambda_d(x)
=\frac12\frac{\mathrm d}{\mathrm dt}
\int_{\R^d}\nabla u^\mathsf TG\nabla u\,\mathrm d\lambda_d(x).
\]
To daje tożsamość energii.

Dla skończonej prędkości propagacji używamy lokalnej gęstości
$e=(|u_t|^2+\nabla u^\mathsf TG\nabla u)/2$
i strumienia $J=-u_tG\nabla u$. Z równania otrzymujemy
$\partial_te+\operatorname{div}J=f u_t$.
Na brzegu ruchomej kuli mamy
$|J\cdot\nu|\le |u_t|\,|G\nabla u\cdot\nu|
\le\sqrt\Lambda |u_t|\sqrt{\nabla u^\mathsf TG\nabla u}
\le\sqrt\Lambda\,e$.
Zatem przy ruchu brzegu z szybkością $-c$ całka energii
w malejącej kuli jest nierosnąca, jeśli $f=0$ wewnątrz.
Zerowe dane $u_1$ i $\nabla u_0$ oznaczają zerową energię
na podstawie. Stąd $u_t=\nabla u=0$ we wnętrzu stożka;
ponieważ $u_0=0$ na podstawie, otrzymujemy tam $u=0$.
\end{proof}

\section{Warunki brzegowe: dlaczego nie można narzucić dowolnych danych?}
Na obszarze ograniczonym energia układu symetrycznego spełnia
prawo~\eqref{eq:pde21-flux}, ale po całkowaniu pojawia się
strumień przez brzeg $\partial\Omega$:
$-\int_{\partial\Omega}u^\mathsf TA_\nu u\,\mathrm d\mathcal H^{d-1}$.
W szczególności dla zagadnienia Cauchy'ego na półprostej
$x>0$ z macierzą diagonalną $A=\operatorname{diag}(a_1,\ldots,a_m)$
warunki brzegowe podaje się tylko na składowych z $a_k>0$,
które wchodzą do obszaru przez $x=0$; dla $a_k<0$
wartości brzegowe są wyznaczane przez wnętrze.

\begin{przyklad}[Dwa kierunki transportu na półprostej]
Dla $a>0$ rozwiązanie $u_t+au_x=0$ na $x>0$ wymaga
jednej wartości brzegowej $u(t,0)=g(t)$: charakterystyki
wychodzące z brzegu wchodzą do obszaru. Natomiast dla $a<0$
charakterystyki biegną ku brzegowi i wartość $u(t,0)$
wynika już z danej początkowej. Dodatkowe narzucenie
niezależnego $g$ może prowadzić do braku rozwiązania.
\end{przyklad}


\begin{twierdzenie}[Dysypatywne dane brzegowe na półprostej]
\label{thm:pde21-halfline-energy}
Niech $A=A^{\mathsf T}\in\R^{m\times m}$ będzie macierzą stałą,
$T>0$, $U\in C^1([0,T]\times[0,\infty);\R^m)$ i
$F\in C([0,T]\times[0,\infty);\R^m)$ spełniają
$U_t+A U_x=F$ na $(0,T)\times(0,\infty)$.
Załóżmy, że $U$, $F$ i ich pochodne mają zanik pozwalający
na całkowanie przez części przy $x\to\infty$ oraz że
$U(t,\cdot),F(t,\cdot)\in L^2(0,\infty)$.
Niech $P_+$ będzie rzutem ortogonalnym na sumę przestrzeni
własnych $A$ odpowiadających dodatnim wartościom własnym,
a $\lambda_+=\max(0,\lambda_{\max}(A))$. Jeżeli
$P_+U(t,0)=g(t)\in\operatorname{Ran}P_+$, to
\[
 \|U(t)\|_{L^2(0,\infty)}\le
 \|U(0)\|_{L^2(0,\infty)}+
 \sqrt{\lambda_+}\,\|g\|_{L^2(0,t)}+
 \int_0^t\|F(r)\|_{L^2(0,\infty)}\,\mathrm dr.
\]
W szczególności dla $g=F=0$ energia nie rośnie.
Nie narzuca się niezależnych danych na ujemnej podprzestrzeni
spektralnej $A$; warunek dla zerowej wartości własnej
wynika z równania zwyczajnego w każdym punkcie przestrzeni.
\end{twierdzenie}
\begin{proof}
Tożsamość lokalna dla macierzy stałej ma postać
$\partial_t|U|^2+\partial_x(U^{\mathsf T}AU)=2U^{\mathsf T}F$.
Po scałkowaniu po półprostej otrzymujemy
\[
 \frac{\mathrm d}{\mathrm dt}\|U(t)\|_2^2
 =U(t,0)^{\mathsf T}AU(t,0)+
 2\langle U(t),F(t)\rangle.
\]
Rozkład spektralny $A=A_+-A_-$, gdzie $A_\pm\ge0$
i ich nośniki spektralne są ortogonalne, daje
$U^{\mathsf T}AU\le\lambda_+|P_+U|^2
=\lambda_+|g|^2$ na $x=0$. Oznaczmy
$E(t)=\|U(t)\|_2^2$, $b(t)=\|F(t)\|_2$ oraz
\[
 a(t)=E(0)+\lambda_+\int_0^t|g(r)|^2\,\mathrm dr.
\]
Z nierówności Cauchy'ego--Schwarza dostajemy
$E'(t)\le a'(t)+2\sqrt{E(t)}\,b(t)$.
Dla $\delta>0$ wprowadzamy
\[
 z_\delta(t)=\left(a(t)+\delta^2+
 2\int_0^t\sqrt{E(r)}b(r)\,\mathrm dr\right)^{1/2}.
\]
Wtedy $\sqrt{E(t)}\le z_\delta(t)$ i
$z_\delta(t)\ge\sqrt{a(t)+\delta^2}$. Dlatego
\[
 z_\delta'(t)=\frac{a'(t)+2\sqrt{E(t)}b(t)}
 {2z_\delta(t)}
 \le\frac{a'(t)}{2\sqrt{a(t)+\delta^2}}+b(t).
\]
Całkując od zera do $t$, wykorzystujemy tożsamość
$\frac{\mathrm d}{\mathrm dt}\sqrt{a(t)+\delta^2}
=a'(t)/(2\sqrt{a(t)+\delta^2})$.
Otrzymujemy
$\sqrt{E(t)}\le\sqrt{a(t)+\delta^2}+
\int_0^t b(r)\,\mathrm dr$. Granica $\delta\downarrow0$
oraz nierówność $\sqrt{r+s}\le\sqrt r+\sqrt s$
kończą dowód. Wartości własne dodatnie odpowiadają
charakterystykom wchodzącym do $x>0$ przez $x=0$,
a ujemne --- wychodzącym, na których nie można niezależnie
zadawać śladu. Dla zerowej wartości własnej projekcja
równania ma postać $(P_0U)_t=P_0F$.
\end{proof}


\section{Podsumowanie}
Symetria części głównej daje tożsamość energii, z której wynikają
jednoznaczność, stabilność i skończona prędkość propagacji.
Po zróżniczkowaniu równania komutatory nie zawierają pochodnych
wyższego rzędu niż $s$; dzięki temu otrzymujemy poprawność
w $H^s$ bez utraty pochodnych. Nie każda macierz o rzeczywistych
wartościach własnych ma te własności; pokazuje to blok Jordana.
W następnym rozdziale porównamy tę metodę z innym
mechanizmem ewolucji: unitarnością dla równania Schrödingera.
Nie jest ono ogólnym układem symetryczno-hiperbolicznym
pierwszego rzędu w przestrzeni, toteż stożkowa zasada
propagacji nie przechodzi na nie bezpośrednio.

